\section{An implementer's guide}
\label{sec:impl}

This section is the advice we wish the papers had included. It assumes you are building on
Calcite (all of \spes{}, \QED{} and \sqlsolver{} consume relational plans, and Calcite's
\texttt{RelNode} is the easiest source).

\subsection{Build order}

Do not start with the fancy parts. In dependency order:

\begin{enumerate}
\item \textbf{Plan $\to$ algebra.} Handle \texttt{TableScan}, \texttt{Filter},
      \texttt{Project}, \texttt{Join}, \texttt{Union(all=true)} only. Represent an expression
      as a sum of terms; a term as (scope, body predicate, list of table applications).
\item \textbf{Normalisation to sum-of-products.} Pure syntax; no solver needed. Test by
      round-tripping and by evaluating both forms on small random instances.
\item \textbf{Pointwise term matching} with scope permutation. At this point you can already
      prove filter/projection/join reordering.
\item \textbf{Bring in SMT for predicates.} This is the single biggest capability jump per
      line of code --- it is the whole difference between \udp{} and its successors.
\item \textbf{Stabilisation} via congruence classes. Second biggest jump.
\item \textbf{Primary keys.} One rewrite. Suddenly a large class of real rules becomes
      provable.
\item \textbf{Squash/\texttt{Distinct}/\texttt{GroupBy}}, then the $\neq0$ Boolean encoding.
\item Only then: NULLs, foreign keys, uninterpreted operators, \liastar{}.
\end{enumerate}

\subsection{Data structures}

\begin{lstlisting}[style=pseudo]
Term   = { scope : List<(TableSym, Var, int power)>,   # Sum over R^k(s)
           body  : Formula,                            # the [P] part, SMT-encodable
           value : Expr }                              # the V part (QOp/HOp applications)
Expr   = List<Term>                                    # a sum of terms
\end{lstlisting}

Keep \texttt{body} in your SMT library's AST from the very beginning. Converting a
home-grown predicate type into SMT later is a rewrite you will do twice.

\subsection{Ten traps, each of which will cost you a day}

\begin{enumerate}
\item \textbf{\lstinline|GROUP BY| needs a squash.} \udp{}'s printed de-sugaring
      (Remark~\ref{rem:groupby}) iterates over every row of $R$, not every distinct key. Use
      \QED{}'s two-factor rule: $\sq{\usum{s}\ind{x=k(s)}Q(s)}\times\ind{y=\alpha(\ldots)}$.
      Get this wrong and multiplicities are silently wrong for duplicated keys.

\item \textbf{The primary-key rule for tables with non-key columns.} Eq.~\eqref{eq:pk8},
      $R\rightsquigarrow\lam{k}{\sq{R(k)}}$, applies only when the key is the \emph{whole}
      schema. In general use Eq.~\eqref{eq:pk7} with the functional-dependency symbol $f_R$.
      Almost every write-up quotes the easy form.

\item \textbf{Foreign-key chasing may not terminate.} \udp{} says so explicitly: a cycle in
      the FK graph loops. Either bound the depth or do what \QED{} does --- emit
      $\forall b,k.\ \sq{S(b,k)}\to\sq{R'(k)}$ as one global assertion and let the solver
      unfold it lazily.

\item \textbf{Over-approximation is one-directional.} In \sqlsolver{}'s non-linear path,
      \textsc{unsat} means equivalent, \textsc{sat} means \emph{nothing}. If you report
      ``not equivalent'' on \textsc{sat} you have an unsound tool.

\item \textbf{Permutation blowup.} $\textsf{termEq}$ enumerates bijections between scopes;
      \QED{} caps at 24 and reports $\sim$98\% of cases match on the identity permutation.
      Try identity first, always.

\item \textbf{Linearisation blowup.} $\mathrm{Lin}(X)$ is the ordered Bell number:
      $1,1,3,13,75,541,4683$. Eager linearisation on a term with five variables already
      generates 541 cases \emph{per table}. Use it only inside the complete fragment.

\item \textbf{Congruence classes only find expressions already present.} Given
      $\ind{x+3=10}$ they will not produce $x=7$. If you need that, call SyGuS.

\item \textbf{Finiteness.} \QED{} assumes every table has finite multiplicity at each row so
      that $\infty$ never appears during term comparison. If you model tables as
      $S\to\Nbar$ you must handle $\infty$ arithmetic; if as $S\to\mathbb{N}$ you must not
      claim completeness over all U-semirings.

\item \textbf{Bag \lstinline|EXCEPT ALL| / \lstinline|INTERSECT ALL| are not free.} They need
      truncated subtraction (monus). \QED{} declines and treats them as unknown operators;
      \sqlsolver{} encodes \lstinline|INTERSECT ALL| as
      $\ind{f_l\le f_r}f_l+\notu{\ind{f_l\le f_r}}f_r$, the pointwise minimum. Pick one
      deliberately.

\item \textbf{Implicit casts.} \QED{}'s failure table lists ``typing'' for 6 Calcite and 26
      CockroachDB pairs. Calcite inserts \texttt{CAST} nodes freely; if you model casts as
      uninterpreted functions, two plans that differ only in an implicit cast will not
      unify. Decide early whether to erase casts (unsound in general) or model them.
\end{enumerate}

\subsection{Testing: use a real database}

The cheapest high-value test harness is differential testing against PostgreSQL, and it
catches translation bugs that no amount of staring will.

\begin{lstlisting}[style=pseudo]
for each candidate rewrite (Q1, Q2):
    1. generate small random instances (0-5 rows, heavy duplicates, some NULLs)
    2. run both queries in Postgres, compare as BAGS:
         (Q1 EXCEPT ALL Q2) must be empty  AND  (Q2 EXCEPT ALL Q1) must be empty
    3. if they differ  -> your prover must NOT say "equivalent"
    4. also evaluate your U-/Q-expression directly on the same instance
       and check it reproduces Postgres' multiplicities row by row
\end{lstlisting}

Step 4 is the one people skip and the one that matters: it validates the \emph{translation},
independently of the decision procedure. Two specific settings to include:

\begin{itemize}
\item \textbf{duplicates in the right-hand table of every join} --- this is what separates a
      correct key model from a wrong one. Our running example produces $(1,60)$ versus two
      copies of $(1,30)$ (\S\ref{sec:pkcounter}); if your translation does not reproduce
      exactly that, your \texttt{GroupBy} or key rule is wrong.
\item \textbf{\lstinline|NULL| in every nullable column} --- three-valued \texttt{AND}/\texttt{OR}
      is easy to get subtly wrong, and \textsf{Null} \texttt{And} \textsf{False} $=$
      \textsf{False} is the canonical check.
\end{itemize}

Use a session-local table to shadow a real one when you want a no-key variant, so you never
mutate the reference schema:

\begin{lstlisting}[style=sql]
create temp table s (a int, b int);        -- hides the real s, no primary key
insert into s values (1,100),(1,100);
\end{lstlisting}

\subsection{Getting plans out of Calcite}

For the record, this is all it takes to obtain the trees in \S\ref{sec:running}:

\begin{lstlisting}[style=pseudo]
SchemaPlus root = Frameworks.createRootSchema(true);
root.add("R", tableOf("X","Y", /*rows*/1e6, /*key*/null));
root.add("S", tableOf("A","B", /*rows*/5e4, /*key*/0));   // column 0 unique
FrameworkConfig cfg = Frameworks.newConfigBuilder().defaultSchema(root).build();
Planner p = Frameworks.getPlanner(cfg);
RelNode rel = p.rel(p.validate(p.parse(sql))).rel;
System.out.println(RelOptUtil.toString(rel, SqlExplainLevel.EXPPLAN_ATTRIBUTES));
\end{lstlisting}

Declare uniqueness through the table's \texttt{Statistic}:
\lstinline[style=pseudo]|Statistics.of(rows, ImmutableList.of(ImmutableBitSet.of(0)))|.
That one bit is what makes Calcite's \texttt{AggregateJoinTransposeRule} produce $Q_2$
instead of the compensated \texttt{SUM $\times$ COUNT} plan of \S\ref{sec:calcite-fires} ---
and it is the same bit your prover needs for the key rewrite. A \texttt{Planner} instance is
single-use; build a fresh one per query.
